\chapter{Earnings}\label{ch:s1:earnings}

A firm reports earnings well above what the market expected, its stock jumps, and then, for weeks, it keeps drifting in the same direction. The
\omterm{def:s1:earnings:pead}{post-earnings-announcement drift} is one of the oldest anomalies in accounting research and one of the most studied; it is also one whose public record
has moved against it. Chordia and co-authors found it concentrated in the most illiquid stocks, where trading costs took 70 to 100\% of its paper profits,
and Martineau found it non-existent in large US stocks since 2006. The synthetic market plants a clean version: 1.2\% of drift per standard deviation of
surprise, spread over sixty days. Traded from the day after the announcement and held for sixty days, it gives a Sharpe ratio of 2.66 after costs; held for
five, $-0.77$; and when the drift loses 70\% of its size, as it did in the synthetic market of Book~7, chapter~13, the sixty-day book's Sharpe ratio over the
last four years falls from 3.24 to 0.75. This chapter is about trading around earnings. The build is \texttt{firm.earnstrat}.

\section{Post-earnings-announcement drift}

\begin{definition}[Event study, event window]\label{def:s1:earnings:event}
An \emph{event study}\index{event study} measures the average abnormal return of securities in \emph{event time}, relative to the date of an event of a
given kind, over an \emph{event window}\index{event window} of days before and after it, to separate the event's effect from the market's.
\end{definition}

\begin{definition}[Post-earnings-announcement drift]\label{def:s1:earnings:pead}
\emph{Post-earnings-announcement drift}\index{post-earnings-announcement drift} is the tendency of a stock's abnormal returns to continue in the direction of
its earnings surprise for weeks after the announcement.
\end{definition}

Ball and Brown (1968) and Bernard and Thomas (1989) are the classic studies; Chan, Jegadeesh and Lakonishok found that past earnings surprises and past
returns each predict large drifts in future returns after controlling for the other, and that analysts' forecasts also respond sluggishly to news: a market
that absorbs information gradually. The synthetic market (Book~7, chapter~5) plants exactly that: on each announcement a standard normal surprise, a jump of
three daily specific volatilities per unit of surprise, and a drift of 1.2\% per unit over the next sixty days. Firms announce four times a year (1.59\% of
stock-days are announcement days). \Cref{fig:s1:earnings:car} is the chapter's \omterm{def:s1:earnings:event}{event study}: for surprises above one standard deviation, the
market-adjusted return is 6.94\% on the announcement day and 7.75\% after sixty days; below minus one, $-6.93\%$ and $-8.56\%$.

\begin{omfigure}
\begin{tikzpicture}
\begin{axis}[width=11cm,height=5.2cm,xlabel={\small days from the announcement},ylabel={\small cumulative abnormal return (\%)},
  tick label style={font=\footnotesize},xmin=-5,xmax=60,ymin=-10,ymax=10,grid=major,grid style={gray!25},
  legend style={font=\scriptsize,at={(0.97,0.5)},anchor=east},table/col sep=comma]
\addplot[omDef,very thick] table[x=day,y=positive_pct]{figdata/strategies-1/07-earnings/car.csv};
\addplot[omThm,very thick] table[x=day,y=negative_pct]{figdata/strategies-1/07-earnings/car.csv};
\addplot[black!50,dashed] coordinates {(-5,0) (60,0)};
\legend{surprise above $+1$ s.d.,surprise below $-1$ s.d.}
\end{axis}
\end{tikzpicture}

\omcaption{\omterm{def:s1:earnings:event}{Event study} of the synthetic market's announcements: mean cumulative market-adjusted return from five days before to sixty days after, for
large positive and large negative surprises. Data: \texttt{s1\_earnings.car}.}\label{fig:s1:earnings:car}
\end{omfigure}

The \omterm{def:s1:earnings:event}{event study} also shows the day after: the positive group gives back 0.33 points, from 6.94\% to 6.61\%. That is the planted one-day reversal of
chapter~2 acting on the announcement's jump, and it decides the first implementation choice below.

\section{Analyst revisions}

\begin{definition}[Earnings-revision strategy]\label{def:s1:earnings:revision}
An \emph{earnings-revision strategy}\index{earnings-revision strategy} buys stocks whose analysts have recently raised their earnings forecasts and sells
stocks whose forecasts have been cut, on the evidence that forecasts, and prices with them, adjust to news gradually.
\end{definition}

The same sluggishness that produces the drift shows in forecasts: Chan, Jegadeesh and Lakonishok found analysts' forecasts responding slowly to past news,
especially for the worst performers. A revision signal is the change in the consensus forecast over a month, scaled by price or by the dispersion of
forecasts (Book~7, chapter~11). It needs a point-in-time history of individual analysts' forecasts, which is a commercial data set; the synthetic market has
no analysts, so the chapter's strategy file on revisions rests on the published record.

\section{Guidance and pre-announcements}

Earnings news does not only arrive on the scheduled date. Firms guide their future earnings, update the guidance between reports, and sometimes
pre-announce a result that will miss or beat expectations, usually in the weeks before the quarter's report. Each of these is an event with its own window,
and the drift question applies to each: does the price move fully on the day, or keep going? The engineering is the same as for announcements (a calendar
of events known point in time and an event-time book), and the data problem is harder: guidance is published in press releases and on calls, not in a
standardised filing, and its timestamp has to be the moment it became public. The strategy file on pre-announcements states that no performance record was
verified for this book.

\section{The announcement-day options angle}

\begin{definition}[Announcement premium]\label{def:s1:earnings:premium}
The \emph{announcement premium}\index{announcement premium} is the higher average return of stocks in the days around their scheduled earnings
announcements, compared with the same stocks at other times.
\end{definition}

Savor and Wilson found that firms scheduled to report earned an annualised abnormal return of 9.9\%, and explained it as a risk premium: announcements
tell investors about other firms and the market too, so the covariance between a firm's news and the market's spikes around them. The synthetic market
plants no \omterm{def:s1:earnings:premium}{announcement premium}, and the portfolio of each day's announcers, bought the close before, earns 1.8 basis points a day over the market with a
$t$ statistic of 0.59: nothing.

What the synthetic market does have is the announcement's jump: the absolute market-adjusted return on an announcement day is 3.02 times that of an
ordinary day. An option expiring just after an announcement has to price that extra move, so its implied volatility should be higher than that of an
option expiring just before, and should fall once the announcement is out. A straddle bought before the announcement is a bet that the move will exceed what the options priced, and a straddle sold is
the opposite bet. Neither is a free lunch: both are trades on the size of the jump relative to its price (Book~9, chapter~4, trades realised against implied volatility around events).

\section{Implementation around event dates}

\begin{center}
\small
\begin{tabular}{@{}lccc@{}}
\toprule
drift book, synthetic market, years 3 to 10 & 5 days & 20 days & 60 days\\
\midrule
bought at the announcement's close: Sharpe ratio after costs & $-3.72$ & 0.02 & 2.03\\
bought a day later: Sharpe ratio before, after costs & 1.38, $-0.77$ & 2.64, 1.51 & 3.31, 2.66\\
bought a day later: net return a year; turnover a year & $-4.7\%$; 65.2 & 4.6\%; 17.2 & 5.5\%; 6.6\\
on the announcement-day return instead of the surprise & $-0.85$ & 0.85 & 2.03\\
drift cut by 70\% from year 7: Sharpe ratio in years 7--10 & $-1.19$ & $-0.51$ & 0.75\\
\bottomrule
\end{tabular}
\end{center}

Three choices decide the result. \emph{When to enter}: buying at the announcement's close means holding through the next day's reversal of the jump, which
turns the five-day book's gross Sharpe ratio negative; waiting a day costs a day of drift and avoids it. \emph{How long to hold}: the drift accrues slowly
(1.2\% per unit of surprise over sixty days, about two basis points a day), so a short holding period pays a full round trip for a few days of it; the
five-day book trades 65 times its gross a year and its costs, 13\% a year, exceed its gross return. \emph{What to measure}: the reported surprise beats the
announcement-day return as a signal (2.66 against 2.03 at sixty days), because the return mixes the surprise with the day's specific noise; real books
combine both, and add revisions.

And the drift can shrink. On the synthetic market whose drift loses 70\% of its size from year 7 (\texttt{MarketConfig.pead\_break}), the sixty-day book
earns 4.8\% a year before the break and 1.4\% after, and its Sharpe ratio falls from 2.18 to 0.75; the twenty-day book goes from 1.48 to $-0.51$. The public
record points the same way: Chordia and co-authors' 0.04\% a month in the most liquid stocks, and Martineau's drift that vanished from large stocks after
2006.

\begin{omfigure}
\begin{tikzpicture}
\begin{axis}[ybar,bar width=10pt,width=10cm,height=5.2cm,xlabel={\small holding period (days)},ylabel={\small Sharpe ratio after costs},
  tick label style={font=\footnotesize},symbolic x coords={5,20,60},xtick=data,ymin=-4.5,ymax=4,grid=major,grid style={gray!25},
  legend style={font=\scriptsize,at={(0.5,1.03)},anchor=south,legend columns=3},table/col sep=comma,enlarge x limits=0.25]
\addplot[fill=omThm!60,draw=omThm] table[x=hold,y=same_day]{figdata/strategies-1/07-earnings/holds.csv};
\addplot[fill=omDef!70,draw=omDef] table[x=hold,y=next_day]{figdata/strategies-1/07-earnings/holds.csv};
\addplot[fill=black!25,draw=black!60] table[x=hold,y=cut_late]{figdata/strategies-1/07-earnings/holds.csv};
\legend{bought at the announcement,bought a day later,a day later; drift cut (years 7--10)}
\end{axis}
\end{tikzpicture}

\omcaption{The drift book's Sharpe ratio after 10 basis points per unit traded, by holding period and entry day, and after the drift loses 70\% of its size.
Data: \texttt{s1\_earnings.run}.}\label{fig:s1:earnings:holds}
\end{omfigure}

\section{Strategy files}

\begin{strategyfile}{SUE drift}\label{strat:s1:earnings:sue}
\sfield{Who pays you, and why} Investors who underreact to earnings news and adjust over weeks.
\sfield{Instruments and venues} Stocks, long positive and short negative surprises.
\sfield{Signal} Standardised unexpected earnings (Book~7, chapter~11) or the surprise against consensus, known at the announcement.
\sfield{Sizing and execution} Event-time book, entered the day after, held up to sixty days.
\sfield{Costs} The strategy's weak point: the drift lives in illiquid stocks where costs are highest.
\sfield{How it dies} Faster price discovery on the announcement day; costs; crowding.
\sfield{Horizon, capacity, infrastructure} Weeks; small capacity in the stocks where the drift remains; an announcement calendar and surprise data.
\sfield{Backtest honestly} Announcement timestamps (before or after the market); consensus as it stood; costs by liquidity.
\sfield{Sources} Chan, Jegadeesh and Lakonishok (1996); Chordia et al.\ (2009): 0.04\% a month in the most liquid stocks, 2.43\% in the most illiquid, costs
70--100\% of paper profits; Martineau (2022): none in large stocks since 2006.
\end{strategyfile}

\begin{strategyfile}{Revision momentum}\label{strat:s1:earnings:revision}
\sfield{Who pays you, and why} As for the drift: analysts and prices adjust gradually.
\sfield{Instruments and venues} Stocks with analyst coverage.
\sfield{Signal} The one-month change in consensus forecasts, scaled by price or dispersion.
\sfield{Sizing and execution} Monthly sorts, or a signal in a multi-factor book.
\sfield{Costs} Moderate turnover.
\sfield{How it dies} As for the drift; herding in forecasts.
\sfield{Horizon, capacity, infrastructure} Months; a point-in-time forecast history.
\sfield{Backtest honestly} Forecasts dated when issued, not when entered into the database; stale forecasts removed.
\sfield{Sources} Chan, Jegadeesh and Lakonishok (1996) on sluggish forecasts; no performance figure verified for this book.
\end{strategyfile}

\begin{strategyfile}{Announcement premium}\label{strat:s1:earnings:premium}
\sfield{Who pays you, and why} Investors who demand a premium for the market-wide risk carried by announcers.
\sfield{Instruments and venues} Stocks with scheduled announcements.
\sfield{Signal} The announcement date, known in advance.
\sfield{Sizing and execution} Long the stocks announcing in the coming days, hedged with the market.
\sfield{Costs} A round trip around each announcement.
\sfield{How it dies} As the premium is a risk premium, it pays with the risk: losses when announcements bring bad market news.
\sfield{Horizon, capacity, infrastructure} Days; an accurate forward calendar.
\sfield{Backtest honestly} Scheduled dates as known in advance, not the dates realised.
\sfield{Sources} Savor and Wilson (2016): an annualised abnormal return of 9.9\% for firms scheduled to report.
\end{strategyfile}

\begin{strategyfile}{Straddle before announcements}\label{strat:s1:earnings:straddle}
\sfield{Who pays you, and why} Whoever sold the options too cheaply, if the move is larger than priced; the buyer pays otherwise.
\sfield{Instruments and venues} Listed equity options expiring just after the announcement.
\sfield{Signal} The announcement move implied by the options against a forecast of the actual move.
\sfield{Sizing and execution} Delta-hedged straddles bought or sold before, closed after.
\sfield{Costs} Option spreads, wide around events.
\sfield{How it dies} Implied moves that price the event correctly on average.
\sfield{Horizon, capacity, infrastructure} Days; option data and a move model.
\sfield{Backtest honestly} Option prices as quoted, bid and ask; the implied event move separated from ordinary volatility.
\sfield{Sources} This chapter's simulation only (announcement moves 3.02 times ordinary ones); no public performance record verified.
\end{strategyfile}

\begin{strategyfile}{Pre-announcement trading}\label{strat:s1:earnings:preannounce}
\sfield{Who pays you, and why} Underreaction to guidance updates and pre-announcements, if it exists.
\sfield{Instruments and venues} Stocks.
\sfield{Signal} Guidance changes and pre-announcements, from press releases, timestamped when public.
\sfield{Sizing and execution} Event-time book as for the drift.
\sfield{Costs} As for the drift.
\sfield{How it dies} As for the drift.
\sfield{Horizon, capacity, infrastructure} Weeks; a news feed parsed for guidance.
\sfield{Backtest honestly} Public timestamps only; no information before release (chapter~17 on news).
\sfield{Sources} No performance record verified for this book.
\end{strategyfile}

\section{Tutorial: sixty days of drift}\label{tut:s1:earnings:tut}

\begin{tutorial}
\textbf{Goal.} Trade the planted drift in \omterm{def:s1:earnings:event}{event time}, choose the entry day and the holding period, and watch the drift shrink. \textbf{End state:} the
table and \cref{fig:s1:earnings:holds}.
\begin{tutsteps}
\item \textbf{The event book}: names with a large recent surprise, from a delay after the announcement, for a holding window.
\omcode{code/firm/earnstrat/firm_earnstrat.py}{34}{48}{The event-time book.}\label{lst:s1:earnings:book}
\item \textbf{The \omterm{def:s1:earnings:event}{event study}}: cumulative abnormal returns in \omterm{def:s1:earnings:event}{event time}.
\omcode{code/firm/earnstrat/firm_earnstrat.py}{57}{67}{Event-time cumulative abnormal returns.}\label{lst:s1:earnings:car}
\item \textbf{Run} \texttt{s1\_earnings.run(hold, signal, cut, delay)} for 5, 20 and 60 days, \texttt{car}, \texttt{moves}, \texttt{announcers} and
\texttt{fig\_earnings.py}.
\end{tutsteps}
\textbf{What to change next.} Weight positions by the size of the surprise; combine the surprise and the announcement return; restrict the book to the 500
most liquid names and measure what the drift is worth there.
\end{tutorial}

\section{Build: earnings strategies}\label{bld:s1:earnings:earnstrat}

\begin{build}
\bfield{Purpose} Event calendars, event-time books and event studies for earnings and other dated events.
\bfield{Interface} \texttt{calendar(events, T, N)}, \texttt{event\_book(flag, surprise, listed, hold, threshold, delay)}, \texttt{announcement\_move(ret,
flag)}, \texttt{event\_car(abret, flag, surprise, before, after, threshold)}.
\bfield{Rules} Events dated when public; positions from the close after the signal; holding windows fixed in advance.
\bfield{Acceptance tests} \texttt{code/firm/earnstrat/tests/}: the calendar and the book by hand, with and without a delay; the move ratio and the event
path on a planted example.
\bfield{Stretch} SUE from point-in-time EPS (Book~7, chapter~11's \texttt{firm.fundpit}); revision signals; the announcement calendar as a forward schedule.
\end{build}

\begin{omsources}
\item R.~Ball and P.~Brown, ``An empirical evaluation of accounting income numbers'', \emph{Journal of Accounting Research} 6(2), 1968.
\item V.~L.~Bernard and J.~K.~Thomas, ``Post-earnings-announcement drift: delayed price response or risk premium?'', \emph{Journal of Accounting Research}
27, 1989.
\item L.~K.~C.~Chan, N.~Jegadeesh and J.~Lakonishok, ``Momentum strategies'', \emph{Journal of Finance} 51(5), 1996.
\item P.~Savor and M.~Wilson, ``Earnings announcements and systematic risk'', \emph{Journal of Finance} 71(1), 2016.
\item T.~Chordia, A.~Goyal, G.~Sadka, R.~Sadka and L.~Shivakumar, ``Liquidity and the post-earnings-announcement drift'', \emph{Financial Analysts
Journal} 65(4), 2009.
\item C.~Martineau, ``Rest in peace post-earnings announcement drift'', \emph{Critical Finance Review} 11(3--4), 2022.
\end{omsources}

\section{Exercises}

\begin{exercise}[$\star$]\label{exo:s1:earnings:1}
With a drift of 1.2\% per unit of surprise over sixty days, what does a surprise of 1.5 standard deviations drift in total, and how much is that a day per
unit of surprise?
\end{exercise}

\begin{exercise}[$\star$]\label{exo:s1:earnings:2}
The five-day book trades 65.2 times its gross a year and the sixty-day book 6.6 times. What does ten basis points per unit traded cost each a year?
\end{exercise}

\begin{exercise}[$\star$]\label{exo:s1:earnings:3}
The synthetic market reverses 5\% of each day's specific shock the next day. What does that take back from an announcement-day move of 6.94\%?
\end{exercise}

\begin{exercise}[$\star\star$]\label{exo:s1:earnings:4}
Why does the reported surprise beat the announcement-day return as a signal here, and when might the return be the better one?
\end{exercise}

\begin{exercise}[$\star\star$]\label{exo:s1:earnings:5}
Why is the announcement date a legitimate signal while the surprise is not known until the announcement? What must a backtest know about the timestamp?
\end{exercise}

\begin{exercise}[$\star\star$]\label{exo:s1:earnings:6}
Announcement-day moves are three times ordinary ones. What does that imply for the implied volatility of an option expiring just after the announcement,
compared with one expiring just before?
\end{exercise}

\begin{exercise}[$\star\star\star$]\label{exo:s1:earnings:7}
\emph{Coding.} Run \texttt{run(60, 'surprise', True)}. Report the Sharpe ratios before and after the break and the annual returns, and say how long you
would need to watch to detect the change (Book~7, chapter~13).
\end{exercise}

\begin{exercise}[$\star\star\star$]\label{exo:s1:earnings:8}
\emph{Find the flaw.} ``Our SUE backtest on all US stocks since 1990 earns 2\% a month; we will run it on the S\&P 500.''
\end{exercise}

\section{Problem: Sixty Days of Drift}

\begin{problem}[{Weekend problem --- an anomaly, traded and fading}]\label{pb:s1:earnings:1}
The planted post-earnings drift of \texttt{firm.synthmkt}, and the public record.

\medskip\noindent\textbf{Part I --- The effect.}
\begin{enumerate}
\item Define an \omterm{def:s1:earnings:event}{event study}, an \omterm{def:s1:earnings:event}{event window} and the drift.
\item What did Chan, Jegadeesh and Lakonishok find?
\item Describe the synthetic market's planted drift and its announcement frequency.
\item Read the \omterm{def:s1:earnings:event}{event study}: announcement-day and sixty-day returns for both groups.
\end{enumerate}

\medskip\noindent\textbf{Part II --- The book.}
\begin{enumerate}[resume]
\item Describe the event-time book.
\item Give the Sharpe ratios at 5, 20 and 60 days, bought at the announcement and a day later.
\item Why does buying at the announcement's close hurt?
\item Why do short holding periods lose after costs?
\end{enumerate}

\medskip\noindent\textbf{Part III --- Variants.}
\begin{enumerate}[resume]
\item Surprise or announcement-day return: which wins, and why?
\item What is the \omterm{def:s1:earnings:premium}{announcement premium}, and what does the synthetic market show?
\item How large are announcement moves, and what does that mean for options?
\item What are revision and pre-announcement strategies, and what data do they need?
\end{enumerate}

\medskip\noindent\textbf{Part IV --- The verdict.}
\begin{enumerate}[resume]
\item State the \emph{named result}: the drift book's Sharpe ratio at holding periods of 5, 20 and 60 days, and its fall after the planted break.
\item What does the public record say about the drift in liquid and large stocks?
\item Where would you still look for it?
\item How would you monitor a live drift book for decay?
\item What are the timestamp traps of earnings data?
\item Which strategy file would you run, and at what size?
\item What does the \omterm{def:s1:earnings:premium}{announcement premium} teach about risk and anomalies?
\item In one sentence: what is an earnings strategy?
\end{enumerate}
\end{problem}

\section{Interview questions}

\begin{interviewq}[$\star$ \iqroles{researcher}]\label{iq:s1:earnings:1}
What is \omterm{def:s1:earnings:pead}{post-earnings-announcement drift}, and why might it exist?
\end{interviewq}

\begin{interviewq}[$\star\star$ \iqroles{researcher, developer}]\label{iq:s1:earnings:2}
How would you build a point-in-time earnings-surprise signal? What can go wrong?
\end{interviewq}

\begin{interviewq}[$\star\star$ \iqroles{researcher}]\label{iq:s1:earnings:3}
An anomaly was strong in the 1990s and is weak today. How do you decide whether to keep trading it?
\end{interviewq}

\begin{interviewq}[$\star\star$ \iqroles{trader}]\label{iq:s1:earnings:4}
How would you trade an earnings announcement with options?
\end{interviewq}

\begin{interviewq}[$\star\star$ \iqroles{risk}]\label{iq:s1:earnings:5}
What risks does an event-driven earnings book carry that a factor book does not?
\end{interviewq}

\begin{interviewq}[$\star\star\star$ \iqroles{researcher}]\label{iq:s1:earnings:6}
A drift of $d$ per unit of surprise accrues evenly over $H$ days; trading costs $c$ per unit traded. Find the holding period that maximises the book's net
return per unit of capital when positions enter and leave once.
\end{interviewq}
